\noindent\null\hfill{\small\textbf{ON THE HANOVER TRIANGULATION.}}\hfill\makebox[0pt][r]{\fontsize{13}{16}\selectfont \textbf{375}}\par\medskip

\begin{center}
\textsc{Gauss} to \textsc{Olbers}.\quad Göttingen, 25~February 1825.
\end{center}

\par\noindent\hspace*{\parindent}\mbox{.\ .\ .\ .\ .\ .}\quad Now a few more words about geodetic levelling.
I am myself \textls[120]{quite clear} on the matter, but I do not know whether I can explain myself so briefly and still bring you to agree with me at once.

\par\noindent\hspace*{\parindent}I have always thought that the expression “local attraction” was very ill chosen and could easily give rise to mistaken views.
One should say that the directions of gravity do not follow the course they would take on a uniform spheroid.
The direction of gravity is the combined result of the attraction of all the constituents of the body of the Earth (and of centrifugal force), and given its irregular composition with respect to density and the unevenness of its surface, that direction cannot be the same as it would be on a regular spheroid.
Yet \textls[120]{whatever the composition may be}, through any point it is always possible to pass a surface that goes all the way around the Earth and to which the direction of gravity is \textls[120]{exactly} perpendicular, and the surface of a continuous liquid at rest would represent it.
\textbf{This} is the surface called a level surface (couche de niveau). To the points of \textls[120]{this} surface \textls[120]{we assign the same height}, without concerning ourselves in the least with whether or how much they deviate from an elliptic spheroid, and the heights above this surface are given both by the barometer and by trigonometric measurement, so that the two must always agree.
Here it is merely assumed\textsuperscript{*}) that the direction of gravity changes along each triangle side according to the law of continuity (though perhaps faster or more slowly than on the elliptic spheroid), and this assumption can give rise to a slight error only if at one of the triangulation stations there is a \textls[120]{genuine} local attraction that is merely local and confined to a small area

\par\smallskip\noindent\rule{.24\textwidth}{.4pt}\par\smallskip
{\footnotesize
\noindent\hspace*{.06\linewidth}\textsuperscript{*}) And of course it is also assumed that all zenith distances are measured reciprocally, as is the case in my survey \textls[120]{without exception}.
For \textls[120]{one-sided} measurements your remark is entirely justified, since the amplitude must then be calculated spheroidally.\par
}

\clearpage
\noindent\makebox[0pt][l]{\fontsize{13}{16}\selectfont \textbf{376}}\hfill{\small\textbf{CORRESPONDENCE WITH OLBERS.}}\hfill\null\par\medskip
\noindent (outside which it is imperceptible).
Yet I am convinced that, with at most the exception of the Brocken, no \textls[120]{such} local attraction occurs anywhere within the whole extent of my triangles and \textsc{Schumacher}'s.\ \mbox{.\ .\ .\ .\ .\ .}
